\chapter{Présentation du contexte}
% TODO: adapter le titre, par exemple « Présentation du contexte du stage » /
% « Présentation du contexte du projet » / « Présentation du contexte du PFE ».

\minitoc

\section{Présentation de l'entreprise}
% Supprimer ou renommer cette section pour un projet universitaire sans entreprise d'accueil.

\subsection{Présentation générale}
% TODO: présenter en 1 court paragraphe l'activité, la taille, la localisation, la création, etc.
[Décrivez l'entreprise d'accueil.]

\subsection{Domaine d'activité}
% TODO: préciser les secteurs d'activité de l'entreprise.
[Décrivez le domaine d'activité.]

\subsection{Objectifs de l'entreprise}
% TODO: présenter la mission et les objectifs de l'entreprise en lien avec votre projet.
[Décrivez les objectifs de l'entreprise.]

\section{Présentation du projet}
% TODO: remplacer le titre par le nom du projet ou de la plateforme.

\subsection{Problématique}
% TODO: à quel problème répond le projet ? Quelles étaient les lacunes ou les difficultés ?
[Décrivez la problématique.]

\subsection{Solution proposée}
% TODO: présenter la solution dans ses grandes lignes et sa réponse à la problématique.
[Décrivez la solution proposée.]

\section{Contexte et objectifs du travail}
% TODO: préciser votre périmètre d'intervention au sein du projet et les objectifs
% concrets qui vous ont été confiés.
[Décrivez le contexte et les objectifs de votre contribution.]

\section{Conclusion}
% TODO: conclure le chapitre en 2 à 3 phrases et introduire le suivant.
[Rédigez une brève conclusion du chapitre.]
